\chapter{Structure and Outline of the Work}
\label{cha:Structure}

The structure, outline, and scope of a scientific work vary depending on the type and topic of the work. However, the following outline scheme has proven to be effective and should generally be adhered to:
%
\begin{itemize}
	\item Cover page and title page
	% \item \emph{Confidentiality clause}
	\item Declaration of authenticity
	\item Task description
	\item Summary and abstract
	\item Table of contents
	\item Introduction
	\item Fundamentals
	\item Main part
	\item Summary and outlook
	\item \textit{List of figures, tables, symbols, and abbreviations}
	\item Bibliography
	\item \textit{Appendices}
\end{itemize}
%
The items in italics are to be considered optional. The following sections briefly discuss each of the outline points.

It is essential to ensure that no more than three levels of organization are used, namely chapters, subchapters, and sub-subchapters. If further levels of structure are required, the basic concept must be reconsidered and modified accordingly.

\section{Title Page}
\label{sec:TitlePage}


The cover page and the title page are the first two pages of this guide or template. They must be designed accordingly and contain all the information listed on them. On the cover page, this includes the logo of the institute, the logo of Leibniz University Hannover, the name of the author, the title of the work, the inventory number of the work (which is assigned shortly before the submission of the work by the institute and is to be placed in the box), and the address of the institute. On the title page, the passage \textit{Preparation for the Module Scientific Work in Engineering}, the course of study, the matriculation number, the submission date, as well as the examiner and the supervisor according to this template, are to be added.

The cover page is to be printed on the outside of the cover of the bound, printed edition. The title page forms the first page of paper that is bound into the cover. Sample copies can be viewed at the Institute for Wind Energy Systems upon request.


%\section{Confidentiality Statement}
%\label{sec:Confidentiality}
%
%For scientific works created in collaboration with a company, a confidentiality statement for the work may be required. The confidentiality statement protects the confidential data of the company from the public (e.g., competitors of the company) and should be included as the first page immediately after the title page in the work. Optionally, the confidentiality statement can be placed directly on the title page.\\
%The content of the confidentiality statement should be coordinated with the supervisor.

\section{Affidavit}
\label{sec:Affidavit}

An affidavit is an important part of a scientific paper and must precede the paper. In the affidavit, the author affirms that this thesis was written independently and without the help of third parties. Furthermore, it is declared that all sources used have been identified as such and listed in the bibliography. The author also confirms that the thesis has not been published in this form in Germany or abroad, nor has it been submitted or submitted for a degree course. In particular, it is also confirmed that artificial intelligence has either not been used to create content or that corresponding passages have been identified. Appendix~\ref{cha:AI-Documentation} shows how the use of AI tools can be documented following the guideline of Leibniz University Hannover.

The wording of the declaration in lieu of an oath must correspond to that on page \pageref{Affidavit}. The declaration must be completed with the place, date and signature of the author.


\section{Task Description}
\label{sec:Task}

The task description provides information about the topic of the thesis, contains the subtasks to be completed and documents the due date and the latest possible submission date as well as the names of the examiner and the supervisor. The assignment is issued by the Institute of Wind Energy Systems. The original should be included after the affidavit. The two pages for the assignment of this template are to be understood as placeholders and are to be replaced by the actual assignment issued. 

%\section{Preface}
%\label{sec:Preface}
%
%The preface is not a fixed component of the thesis and is written depending on the theme and type of the work. It precedes the abstract of the work and therefore does not contain any substantive aspects of the work.\\
%In general, a preface provides background information on the creation of the work from the author's perspective. This may include personal reflections, insights, and challenges encountered during data collection, as well as the author's motivation and unique perspectives on addressing the theme. Acknowledgments to sponsors or supporters can also be included in the preface.\\
%
%The length of the preface should not exceed one DIN A4 page.


\section{Abstract and "Kurzfassung"}
\label{sec:Abstract_EN}

A maximum of half a page should be devoted to an abstract of the paper. The abstract should give the reader an overview of the content of the document. It should therefore contain a brief introduction to the topic, the methodological approaches used to work on the topic and the main results of the work. An abstract in English should be added after the abstract in German. Both abstracts should not exceed one page in total.

\section{Table of Contents}
\label{sec:TableOfContents}

The table of contents includes all chapters, sub-chapters and sub-sub-chapters in the document. Other lists (lists of figures, tables, symbols and abbreviations) and any appendices must also be included in the bibliography. 

The headings in the text and in the table of contents must match exactly. Headings of individual bullet points must not coincide with the overall topic of the paper or with the headings of sub-items. Care must be taken to ensure that the expected content is clear from the headings. Generally valid chapter titles such as \textit{Fundamentals} or \textit{Main part} should be avoided.

The page number of the beginning of the section must be listed next to the numbered section titles.

\section{Introduction}
\label{sec:Introduction}

The introduction is the first numbered chapter of the paper. It introduces the topic, motivates the work and derives the problem. The introduction serves to place the topic in an overarching context and, if necessary, to delimit the topic. The introduction concludes with a description of the structure of the paper.



\section{Technical-physical Basics}
\label{sec:Fundamentals}

The second chapter of a scientific paper is dedicated to the technical-physical basics, which a technically educated reader must know in order to be able to understand all the following chapters. Care should be taken to give this chapter a content-related title. A general title such as \textit{Technical-physical basics} is not permitted.



\section{Main Part}
\label{sec:MainPart}

In the main part of a scientific paper, the entirety of the problem being worked on is presented. The methodological approaches and the solution path should be presented, as well as the results of the investigations. 

In general, attention should be paid to a clear presentation of the train of thought, a comprehensible structure and logical arrangement of the content as well as precise, comprehensible, short and concise formulation.

The structure of the content and thus the order of presentation must be based on factual aspects and not on the chronological sequence of the work. The main part may comprise several thematically separate chapters. Care should be taken to give all chapters content-related titles. A generic title such as \textit{Main part} is not permitted.

It must be possible to reproduce the results obtained by describing the investigations carried out. If results or solutions from other sources are cited, these should only be explained to the extent necessary for understanding. A derivation may be omitted. However, derivations of methods actually used should be included in any case.



\section{Summary and Outlook}
\label{sec:Summary}

A scientific paper is concluded with a summary. The summary must contain clear statements about the methods used, the results obtained and their significance, the limits of validity and the progress made compared to the state of knowledge at the beginning of the work and, if applicable, the applicability of the results. 

In addition, the summary should make appropriate reference to the contribution of the work and possibly provide an outlook on possible in-depth research or further investigations of the topic.

A reader who is looking for information on a specific topic should have a good overview of the entire thesis after reviewing the abstract, the summary and the introduction. For this reason, the summary usually repeats statements that have already been described in detail in the main body of the paper.



\section{List of Figures, Tables, Symbols and Abbreviations}
\label{sec:Fig-Tab-Symb-Abk-v}

The purpose of the list of figures is to list all figures and their captions with page numbers. The list of figures is optional, but must be included if the document contains more than three figures. The list of tables must be included in the same way as the list of figures if there are more than three tables in the document.

The list of tables has the same purpose as the list of figures, and is also optional.

The list of symbols is used to list all the symbols used. This is particularly helpful when working with many formulas, as formula symbols are interpreted as symbols. The list is in tabular form and should contain information on the symbols used, their meaning and unit. The list of symbols should be sorted alphabetically. Greek and Roman symbols should be listed in separate blocks.

The list of abbreviations is added to a scientific paper if many unfamiliar abbreviations are used in the body text. Both the abbreviations and their meanings are listed in tabular form. Only abbreviations not listed in the "`Duden"' should be included. 

\section{Bibliography}
\label{sec:Literature}

All cited sources are listed in the bibliography. This must be done in alphabetical order of the first author's surname. The numbers of the sources must correspond to the reference numbers used in the text for the corresponding literature or source.

Further information on citation is provided in section \ref{sec:Citation}. Layout examples for different types of sources can be found in the bibliography at the end of the document. The style of the bibliography should be adopted in all cases.

\section{Appendices}
\label{sec:Appendices}

Appendices in a scientific paper are used to present content such as figures, tables, derivations, program source texts, etc., which do not fit in the text or are unnecessary but are nevertheless important for the comprehensibility and reproducibility of the work. Appendices are listed in the table of contents or in a separate index. Reference should be made to the appendices in the body text.

